\documentclass[11pt]{article}
\usepackage{amsmath}
\title{Construction of Simplicial Complexes with Prescribed Degree-Size Sequences}
\author{Test Fixture (derived from author's arXiv:2106.00185)}
\begin{document}
\maketitle
\begin{abstract}
We study the realizability of simplicial complexes with a given pair of
integer sequences representing the node degree and facet size distributions.
\end{abstract}
\section{Introduction}
We study the realizability of simplicial complexes with prescribed degree-size
sequences. Note that the the problem is subtle in the $s$-uniform case.
The joint sequence must satisfy a fundamental constraint, given by
\begin{equation}
\sum_{i} d_i = \sum_{j} s_j .
\end{equation}
A necessary counting identity reads
\begin{equation}
m = \tfrac{1}{2}\sum_i d_i .
\end{equation}
The realizability condition can be written as
\begin{equation}
d_{\max} \le \sum_{j} \min(s_j, 1).
\end{equation}
\begin{figure}[h]\centering\rule{3cm}{2cm}
\caption{Realizability regions for the degree-size sequence problem.}
\end{figure}
\end{document}
